\section{Design}\label{sec:design} % (fold)
\subsection{Top-level System Architecture}
\begin{figure}[H]
	\centering
	\resizebox{\textwidth}{!}{%
		\begin{tikzpicture}[
			>={Latex[scale=1.2]},
			block/.style={draw=black, thick, minimum width=2.4cm, minimum height=1.5cm, align=center, font=\sffamily\small},
			ctrlblock/.style={draw=black, thick, minimum width=14cm, minimum height=1.2cm, align=center, font=\sffamily\small\bfseries},
			bus/.style={thick, ->},
			bidir/.style={thick, <->},
			ctrlsignal/.style={dashed, ->, font=\sffamily\scriptsize},
			mainbus/.style={line width=3pt, draw=black!60}
			]

			\node[ctrlblock] (ctrl) at (0, 5) {Controller Unit};

			\node[block] (pc)  at (-5, 2)  {Program\\Counter (PC)};
			\node[block] (ir)  at (0, 2)   {Instruction\\Register (IR)};

			\node[block] (mux) at (-5, -1) {Address\\MUX};
			\node[block] (mem) at (-5, -4) {Memory\\($32 \times 8$)};

			\node[block] (alu) at (3.5, -1) {ALU};
			\node[block] (ac)  at (3.5, -4) {Accumulator\\(AC)};

			\draw[mainbus] (-2.0, 0.5) -- (5.5, 0.5) node[pos=0.975, above=4pt, font=\sffamily\scriptsize, text=black] {8-bit Data Bus};

			\draw[bus] (-5.5, 1.25) -- node[above=2pt, pos=0.5, sloped] {\scriptsize{pc\_addr}} (-5.5, -0.25);

			\draw[bus] (-0.6, 1.25) -- (-0.6, 1.0) -- node[above] {ir\_addr[4:0]} (-4.5, 1.0) -- (-4.5, -0.25);

			\draw[bus] (-5, -1.75) -- node[left] {addr[4:0]} (-5, -3.25);
			\draw[bidir] (-3.8, -4) -- node[below] {data} (-2, -4) -- (-2, 0.5);

			\draw[bus] (0, 0.5) -- node[right=2pt] {data\_in} (0, 1.25);
			\draw[bus] (1.2, 2) -- node[above] {opcode} (4, 2) -- (4, -0.25);
			\draw[bus] (3, 0.5) -- node[left=2pt] {inB} (3, -0.25);

			\draw[bus] (3.5, -1.75) -- node[left=2pt] {alu\_out} (3.5, -3.25);
			\draw[bus] (2.3, -4) -- (1.5, -4) -- node[left=2pt] {inA} (1.5, -1) -- (2.3, -1);
			\draw[bus] (4.7, -3.8) -- (5.5, -3.8) -- node[right=2pt, pos=0.2, align=left, sloped] {ac\_out\\(STO)} (5.5, 0.5);

			\draw[ctrlsignal] (-5, 4.4) -- node[right=2pt] {inc\_pc, ld\_pc} (-5, 2.75);
			\draw[ctrlsignal] (-0.5, 4.4) -- node[left=2pt] {ld\_ir} (-0.5, 2.75);

			\draw[ctrlsignal] (-6.5, 4.4) -- node[right=2pt, pos=0.75] {sel} (-6.5, -1) -- (-6.2, -1);
			\draw[ctrlsignal] (-7, 4.4) -- node[left=2pt, pos=0.2] {rd, wr} (-7, -4) -- (-6.2, -4);
			\draw[ctrlsignal] (6.5, 4.4) -- node[right=2pt, pos=0.4] {ld\_ac} (6.5, -4.2) -- (4.7, -4.2);

			\draw[bus] (0.5, 2.75) -- node[right=2pt] {opcode} (0.5, 4.4);

			\draw[bus] (4.7, -1) -- (7.7, -1) -- node[right=2pt] {zero} (7.7, 5) -- (ctrl.east);

			\fill[black!60] (-2, 0.5)  circle (3pt);
			\fill[black!60] (0, 0.5)   circle (3pt);
			\fill[black!60] (3, 0.5)   circle (3pt);
			\fill[black!60] (5.5, 0.5) circle (3pt);

		\end{tikzpicture}
	}
	\caption{Top-Level Architecture of the RISC Processor}
	\label{fig:top_level_arch}
\end{figure}

\subsection{Functional Blocks}
\subsubsection{Program Counter (PC)}
\begin{itemize}
	\item \textbf{Function:}
	      \begin{itemize}
		      \item The counter is a crucial component used to track and count program instructions.
		      \item It can also be utilized to count program states.
		      \item The counter operates synchronously on the rising edge of the \texttt{clk} signal.
		      \item The \texttt{reset} signal is active-high, clearing the counter back to 0.
		      \item The counter has a width of 5 bits.
		      \item It features a load function to force an arbitrary address into the counter. If not loading, the counter increments normally.
	      \end{itemize}
	\item \textbf{Inputs:} \texttt{clk}, \texttt{rst}, \texttt{ld\_pc}, \texttt{inc\_pc}, \texttt{data\_in}.
	\item \textbf{Outputs:} \texttt{pc\_out}.
\end{itemize}

\subsubsection{Address Multiplexer (Address MUX)}
\begin{itemize}
	\item \textbf{Function:}
	      \begin{itemize}
		      \item The Address MUX selects between the instruction address (during the instruction fetch phase) and the operand address (during the instruction execution phase).
		      \item The MUX has a default data width of 5 bits.
		      \item The width is parameterized to allow for easy modification if necessary.
	      \end{itemize}
	\item \textbf{Inputs:} \texttt{pc\_addr}, \texttt{ir\_addr}, \texttt{sel}.
	\item \textbf{Outputs:} \texttt{addr\_out}.
\end{itemize}

\subsubsection{Arithmetic Logic Unit (ALU)}
\begin{itemize}
	\item \textbf{Function:}
	      \begin{itemize}
		      \item The ALU executes arithmetic and logical operations. The specific operation performed depends on the instruction's opcode.
		      \item The ALU executes 8 distinct operations on 8-bit operands (\texttt{inA} and \texttt{inB}).
		      \item It produces a 8-bit output and a 1-bit asynchronous \texttt{zero} flag.
		      \item The \texttt{zero} flag indicates whether the input \texttt{inA} is equal to 0.
	      \end{itemize}
	\item \textbf{Inputs:} \texttt{inA}, \texttt{inB}, \texttt{opcode}.
	\item \textbf{Outputs:} \texttt{alu\_out}, \texttt{zero}.
\end{itemize}

\begin{table}[H]
	\centering
	\caption{ALU Opcode Lookup Table}
	\renewcommand{\arraystretch}{1.3}
	\begin{tabular}{|c|c|p{8.5cm}|c|}
		\hline
		\textbf{Opcode} & \textbf{Code} & \textbf{Operation Description}                                                                                               & \textbf{Output}      \\ \hline
		\texttt{HLT}    & 000           & Halts program execution.                                                                                                     & \texttt{inA}         \\ \hline
		\texttt{SKZ}    & 001           & Checks if the ALU result is equal to 0. If zero, it skips the next instruction; otherwise, execution continues normally.     & \texttt{inA}         \\ \hline
		\texttt{ADD}    & 010           & Adds the value in the Accumulator to the memory value at the instruction's address, returning the result to the Accumulator. & \texttt{inA + inB}   \\ \hline
		\texttt{AND}    & 011           & Performs a bitwise AND on the Accumulator value and the memory value at the instruction's address.                           & \texttt{inA AND inB} \\ \hline
		\texttt{XOR}    & 100           & Performs a bitwise XOR on the Accumulator value and the memory value at the instruction's address.                           & \texttt{inA XOR inB} \\ \hline
		\texttt{LDA}    & 101           & Reads the value from the address in the instruction and loads it into the Accumulator.                                       & \texttt{inB}         \\ \hline
		\texttt{STO}    & 110           & Writes the Accumulator's data to the memory address specified in the instruction.                                            & \texttt{inA}         \\ \hline
		\texttt{JMP}    & 111           & Unconditional jump; jumps to the target address in the instruction and continues execution.                                  & \texttt{inA}         \\ \hline
	\end{tabular}
\end{table}

\subsubsection{Controller Unit}
\begin{itemize}
	\item \textbf{Function:} The central finite state machine that decodes instructions and dictates the operation of all other blocks via control signals.
	\item \textbf{Inputs:} \texttt{clk}, \texttt{rst}, \texttt{opcode}, \texttt{zero}.
	\item \textbf{Outputs:} \texttt{sel}, \texttt{rd}, \texttt{ld\_ir}, \texttt{halt}, \texttt{inc\_pc}, \texttt{ld\_ac}, \texttt{ld\_pc}, \texttt{wr}, \texttt{data\_e}.
\end{itemize}

\begin{table}[H]
	\centering
	\caption{Control Execution Scenarios by Opcode}
	\renewcommand{\arraystretch}{1.3}
	\begin{tabular}{|l|l|l|}
		\hline
		\textbf{Instruction Scenario} & \textbf{Active States}               & \textbf{Expected Control Signals}                              \\ \hline
		\texttt{HLT} (000)            & \texttt{OP\_ADDR}                    & \texttt{halt = 1}                                              \\ \hline
		\texttt{ADD, AND, XOR, LDA}   & \texttt{OP\_FETCH} to \texttt{STORE} & \texttt{rd = 1}, \texttt{ld\_ac = 1} (at \texttt{STORE} phase) \\ \hline
		\texttt{STO} (110)            & \texttt{ALU\_OP}, \texttt{STORE}     & \texttt{wr = 1}, \texttt{data\_e = 1}                          \\ \hline
		\texttt{JMP} (111)            & \texttt{ALU\_OP}, \texttt{STORE}     & \texttt{ld\_pc = 1}                                            \\ \hline
	\end{tabular}
\end{table}

\subsubsection{Instruction Register (IR)}
\begin{itemize}
	\item \textbf{Function:} As described in the reference documentation, it captures and holds the current instruction fetched from memory, splitting it into an opcode and an operand.
	\item \textbf{Inputs:} \texttt{clk}, \texttt{rst}, \texttt{ld\_ir}, \texttt{data\_in}.
	\item \textbf{Outputs:} \texttt{opcode}, \texttt{operand}.
\end{itemize}

\subsubsection{Accumulator Register (AC)}
\begin{itemize}
	\item \textbf{Function:} As described in the reference documentation, it acts as the primary working register for ALU operations and data transfers.
	\item \textbf{Inputs:} \texttt{clk}, \texttt{rst}, \texttt{ld\_ac}, \texttt{data\_in}.
	\item \textbf{Outputs:} \texttt{ac\_out}.
\end{itemize}

\subsubsection{Memory}
\begin{itemize}
	\item \textbf{Function:}
	      \begin{itemize}
		      \item The memory module stores both instructions and data.
		      \item It is designed to separate read and write functions while utilizing a single bidirectional data port. Reading and writing cannot occur simultaneously.
		      \item Configured for 5-bit addresses and 32-bit data.
		      \item Uses 1-bit control signals to enable reading (\texttt{rd}) and writing (\texttt{wr}).
		      \item Memory operations are synchronized to the rising edge of the \texttt{clk} signal.
	      \end{itemize}
	\item \textbf{Inputs:} \texttt{clk}, \texttt{rd}, \texttt{wr}, \texttt{addr}.
	\item \textbf{Inout:} \texttt{data}.
\end{itemize}
% section Design (end)
